% =============================================================================
%  TopoBrain — Manuscript for IET Image Processing
%  Target journal: IET Image Processing  (ISSN 1751-9667; Wiley/IET)
%  Author guidelines: https://ietresearch.onlinelibrary.wiley.com/hub/journal/17519667/homepage/author-guidelines
%
%  Notes for the submitting author:
%    * This file is written against the standard `article` class with twocolumn
%      mode so that it compiles offline (e.g., on Overleaf "free") for
%      pre-submission review.
%    * For the actual Wiley/IET submission, the publisher requires their
%      `wileyNJD-v2.cls` template (downloadable from authorservices.wiley.com
%      after manuscript submission). Drop-in replacement: copy the body of
%      this file (everything between \begin{document} ... \end{document}) into
%      Wiley's main.tex skeleton, fix the \author blocks to their format, keep
%      everything else as-is.
%    * Reference style is numeric (Vancouver-like), matching IET house style.
% =============================================================================

\documentclass[10pt,twocolumn,a4paper]{article}

% ---- Packages -------------------------------------------------------------
\usepackage[utf8]{inputenc}
\usepackage[T1]{fontenc}
\usepackage[margin=20mm,top=20mm,bottom=22mm]{geometry}
\usepackage{graphicx}
\usepackage{amsmath,amssymb}
\usepackage{booktabs}
\usepackage{multirow}
\usepackage{caption}
\usepackage[hyphens]{url}
\usepackage[hidelinks]{hyperref}
\usepackage{cite}
\usepackage{authblk}
\usepackage{xcolor}
\usepackage{abstract}

\graphicspath{{figures/}}
\setlength{\columnsep}{6mm}

% IET-style section heading sizes
\usepackage{titlesec}
\titleformat{\section}{\large\bfseries}{\thesection.}{0.6em}{}
\titleformat{\subsection}{\normalsize\bfseries}{\thesubsection.}{0.5em}{}

% Caption styling (IET house style: bold "Fig. X." / "Table X." prefix)
\captionsetup[figure]{labelfont=bf,labelsep=period,name=Fig.}
\captionsetup[table]{labelfont=bf,labelsep=period}

% =============================================================================
%  TITLE
% =============================================================================

\title{\bfseries TopoBrain: Anatomy-Aware Conditional Diffusion for 3T-to-7T MRI Synthesis with Application to Alzheimer's Disease Biomarkers}

\author[1]{Prabesh Bashyal}
\author[1]{Pratik Adhikari}
\author[1]{Spandan Bhandari}
\author[1,*]{Nanda Bikram Adhikari}

\affil[1]{Department of Electronics and Computer Engineering, Pulchowk Campus, Institute of Engineering, Tribhuvan University, Lalitpur, Nepal}
\affil[*]{Corresponding author: \texttt{adhikari@ioe.edu.np}}
\date{}

\begin{document}
\twocolumn[
  \begin{@twocolumnfalse}
    \maketitle
    \vspace{-1.4em}
    \begin{abstract}
\noindent 7T magnetic resonance imaging offers superior tissue contrast and
spatial resolution that benefits Alzheimer's disease biomarker assessment, but
the scarcity and cost of 7T scanners restricts clinical access. We present
\textbf{TopoBrain}, an anatomy-aware conditional diffusion framework that
synthesises 7T-quality T1-weighted volumes from routine 3T inputs while
preserving the subcortical anatomy required for downstream Alzheimer's
analysis. The architecture is a 22.4-million-parameter 3D U-Net with two
decoders --- a noise-prediction head trained with a denoising diffusion
probabilistic model objective and a 4-class tissue-segmentation head trained
with auxiliary cross-entropy supervision --- coupled with an edge-aware
topology loss that explicitly regularises tissue boundaries. Training on a
paired 3T/7T cohort (10 subjects, 20 T1w volumes) achieves SSIM\,$=$\,0.8991
and PSNR\,$=$\,20.00\,dB on held-out evaluation, with brain-mask
Dice\,$=$\,0.9436 and HD95\,$=$\,2.05\,mm. Ablation against a no-topology
variant reduces whole-brain HD95 by 6.5$\times$ (13.34\,mm to 2.05\,mm) and
suppresses connected-component fragmentation across all tissue classes.
Subcortical regions of interest relevant to early Alzheimer's pathology
--- bilateral hippocampus, amygdala, thalamus and caudate --- are preserved
as essentially single connected components (largest-CC fraction
$\geq$\,99.8\,\%, GM coverage $\geq$\,81.8\,\%). Out-of-distribution
deployment on 30 ADNI baseline 3T scans (15 AD, 15 CN) produces
anatomically coherent outputs, recovers six of six expected tissue-volume
biomarker directions, and yields a CSF/GM atrophy index with Cohen's
$d=+0.681$. We present an end-to-end pipeline (synthesis, parcellation,
classifier-comparison, volumetry-agreement) that is fully implemented and
validated against synthetic data; full-cohort empirical AD classification is
identified as the immediate next study, requiring large-scale
FastSurfer-based parcellation compute beyond the present scope.
    \end{abstract}
    \noindent\textbf{Keywords:} medical image synthesis; diffusion
models; 3T-to-7T MRI; Alzheimer's disease; anatomy-aware deep learning;
neuroimaging biomarkers.

    \vspace{0.8em}
    \hrule
    \vspace{0.6em}
  \end{@twocolumnfalse}
]

% =============================================================================
\section{Introduction}
% =============================================================================

7-Tesla magnetic resonance imaging (MRI) provides higher signal-to-noise
ratio and finer cortical-subcortical contrast than the clinical
standard 3T systems, enabling more sensitive measurement of structural
biomarkers used in Alzheimer's disease (AD) research --- hippocampal
sub-field volumetry, entorhinal cortex thickness, ventricular morphometry,
and cortical layer profiling
\cite{wisse2014hippo,deflores2020hippocampal,yushkevich2015subfields}. In
practice, however, 7T scanners are rare, expensive and concentrated at
research centres; routine clinical AD imaging worldwide is performed at
3T, where partial-volume effects and reduced contrast limit early
detection
\cite{frisoni2010clinical,jack2018niaaa}.

A model that produces 7T-equivalent contrast from a single 3T input, on
the patient's own voxel grid, would let clinical 3T sites access
biomarker pipelines previously confined to 7T research cohorts. The
challenge is that synthesised volumes must remain anatomically faithful:
hallucinated cortex or fragmented subcortical structures would corrupt
downstream parcellation and obscure AD-relevant atrophy. Pixel-similarity
losses (L1, MSE) optimise PSNR but tend to over-smooth high-frequency
detail, and adversarial losses can introduce realistic-looking but
anatomically unfaithful texture
\cite{blau2018perception,isola2017image,zhu2017unpaired}.

We introduce \textbf{TopoBrain}, an anatomy-aware conditional diffusion
framework whose training objective explicitly couples three signals:
(i) the standard denoising diffusion probabilistic model (DDPM) noise-
prediction loss for high-fidelity synthesis
\cite{ho2020denoising}, (ii) an auxiliary 4-class tissue-segmentation
loss supervised by FreeSurfer-derived ground truth from the paired 7T
target, and (iii) an edge-aware topology loss that combines per-class
cross-entropy with Sobel-gradient boundary supervision to suppress
spurious connected components. The architecture (\S\ref{sec:arch}) is a
22.4M-parameter 3D U-Net with dual decoders --- a noise-prediction head
and a tissue-segmentation head --- trained jointly. At inference, both
heads run in 3.1\,min on an NVIDIA A100 for a single full volume.

The contributions of this paper are:

\begin{enumerate}\itemsep1pt
\item An end-to-end conditional-diffusion framework for 3T$\rightarrow$7T
T1w synthesis with explicit segmentation supervision and an edge-aware
topology regulariser, achieving SSIM\,$=$\,0.8991 and brain-mask
HD95\,$=$\,2.05\,mm on held-out paired evaluation.
\item An ablation demonstrating that the topology loss reduces
whole-brain HD95 by 6.5$\times$ and suppresses tissue-component
fragmentation across CSF, grey matter (GM), and white matter (WM).
\item A subcortical region-preservation analysis showing largest-CC
fraction $\geq$\,99.8\,\% across the bilateral hippocampi, amygdalae,
thalami and caudates --- the structures most affected in early AD.
\item An out-of-distribution deployment on 30 ADNI baseline 3T scans
(15 AD, 15 CN), producing anatomically coherent outputs whose tissue
volumes recover six of six expected AD biomarker directions and a
CSF/GM atrophy index with Cohen's $d=+0.681$.
\item A complete, self-validated downstream pipeline (cohort assembly,
synthesis, parcellation, classifier comparison, volumetry agreement)
whose full-cohort empirical evaluation is left as the natural follow-up
study and is reproducible from the openly published code.
\end{enumerate}

% =============================================================================
\section{Related work}
\label{sec:related}
% =============================================================================

\textbf{Cross-field-strength MRI synthesis.} Early CNN-based 3T-to-7T
synthesis (e.g., WATNet \cite{qu2020synthesized}) used domain-specific
spatial and wavelet features but operated patchwise without
anatomical regularisation. Recent hybrid CNN-transformer designs such as
VCT \cite{eidex2024highresolution} report ADC-map synthesis at PSNR
27.0\,dB, SSIM 0.945 on a single dataset, and large-cohort 3T$\rightarrow$
7T transformer models \cite{gicquel2025converting} achieve plausible
contrast translation but do not enforce tissue-class fidelity.
BrainMRDiff \cite{bhattacharya2025brainmrdiff} couples diffusion with
GAN-style adversarial supervision for tumour synthesis, while latent
diffusion has been applied to 1.5T$\rightarrow$3T super-resolution with
downstream AD/MCI classification \cite{yoon2024latent}. None of these
prior systems enforces a joint segmentation head supervised by paired
7T-derived tissue labels, and few report subcortical region-preservation
metrics relevant to AD.

\textbf{Anatomy-aware generative models.} Topology-preserving losses
have appeared in segmentation \cite{hu2019topology,clough2020topo} and
graph super-resolution \cite{singh2024strongly}, with clDice
\cite{shit2021cldice} introduced for tubular vessel structures. We
adapt the underlying intuition --- penalise fragmented connected
components and ill-defined boundaries --- to tissue-class supervision
in a 3D diffusion U-Net, while keeping the loss differentiable and
inexpensive (Sobel-edge cross-entropy and Dice).

\textbf{Alzheimer's disease MRI biomarkers.} Structural-MRI biomarkers
for AD are well established: bilateral hippocampal atrophy
\cite{frisoni2010clinical,jack2018niaaa}, ventricular enlargement,
entorhinal-cortex thinning \cite{wisse2014hippo}, and the CSF/GM ratio
as an atrophy index \cite{frisoni2010clinical}. Automated AD
classifiers from 3T sMRI achieve area-under-curve (AUC) 0.75--0.85 in
published benchmarks \cite{bron2015caddementia,cuingnet2011automatic,
basaia2019automated,liu2022generalizable,islam2018brain}. The relevant
question for synthesis-augmented pipelines is whether 7T-supervised
features extracted from synthesised volumes preserve the AD signal
already accessible at 3T, and whether they enable downstream
hippocampal sub-field measurement that 3T alone cannot.

% =============================================================================
\section{Materials and Methods}
\label{sec:methods}
% =============================================================================

\subsection{Datasets}
\label{sec:data}

\textbf{Paired 3T/7T training cohort.} Ten healthy adult subjects with
two scanning sessions each (3T at session 1 (ses-1), 7T at session 2
(ses-2)) provided 20 T1-weighted volumes. Voxel spacing is approximately
1.0\,mm isotropic, with typical dimensions $256{\times}304{\times}308$.
Subject-level partitioning prevents leakage: nine subjects are used for
training and subject sub-06 is held out for evaluation. The 7T sessions
include FreeSurfer aparc+aseg parcellations used to derive the
4-class tissue ground truth (background, cerebrospinal fluid (CSF),
grey matter, white matter) for the auxiliary segmentation head. Data is
organised in the BIDS specification \cite{gorgolewski2016brain}.

All 7T (training) and 3T (training and ADNI) volumes are preprocessed
identically: HD-BET CNN-based skull stripping
\cite{wegmayr2022hdbet}, N4 bias-field correction, and intensity
normalisation to the diffusion-loss range $[-1, 1]$ via
0.5/99.5-percentile clipping.

\textbf{ADNI external cohort.} Baseline 3T MPRAGE scans were assembled
from the Alzheimer's Disease Neuroimaging Initiative
\cite{petersen2010adni}. The cohort manifest is constructed from four
ADNI study tables (DXSUM diagnosis-of-record, PTDEMOG demographics,
MRI3META MRI conduct flags, MRIQC scan QC) and filtered to baseline
visits with confirmed AD or cognitively normal (CN) diagnosis. After
demographic filtering 266 subjects remain (107\,AD, 159\,CN); 1:1
age/sex matching produces 103\,AD$\leftrightarrow$103\,CN. The present
work reports an external-cohort deployment on 30 subjects (15\,AD,
15\,CN), age-matched (AD\,76.4\,$\pm$\,4.7\,yr, CN\,74.6\,$\pm$\,2.8\,yr).
ADNI does not provide paired 7T volumes; evaluation is therefore
indirect and based on tissue-volume biomarker direction.

\subsection{Architecture}
\label{sec:arch}

The denoising network is a 3D U-Net with explicit 3T conditioning at
each encoder level. The input to the noise-prediction head is the
concatenation of the noisy 7T target $x_t^{7T}$ at diffusion timestep
$t$ and the corresponding registered 3T volume $x^{3T}$ (two channels).
Encoder levels at $\{32,64,128,256\}$ feature widths each contain two
residual blocks with GroupNorm-8 normalisation and SiLU activation; the
$1/8$-resolution bottleneck includes a multi-head 3D self-attention
module over the spatial dimensions. The decoder splits into two heads:

\begin{itemize}\itemsep1pt
\item \emph{Decoder A (noise prediction)} outputs a single-channel
volume $\hat\epsilon_\theta(x_t^{7T},t,x^{3T})$ matching the
DDPM training objective.
\item \emph{Decoder B (tissue segmentation)} outputs a four-channel
volume of class logits over $\{\text{BG},\text{CSF},\text{GM},\text{WM}\}$.
\end{itemize}

The total parameter count is 22.4\,M; the inference memory footprint at
batch size 1 is approximately 3.1\,GB. Time embeddings use sinusoidal
positional encoding mapped through a two-layer MLP into each residual
block. The architecture builds on \cite{ronneberger2015unet,
ho2020denoising,vaswani2017attention} but introduces the dual-decoder
joint training that is central to the anatomy-preservation claim.

\subsection{Training objective}
\label{sec:loss}

The training loss is a weighted sum of three components:
%
\begin{equation}
\mathcal{L}_\text{total} = 1.0\,\mathcal{L}_\text{pixel} + 0.2\,\mathcal{L}_\text{topo} + 0.25\,\mathcal{L}_\text{percep},
\end{equation}
%
where $\mathcal{L}_\text{pixel}$ is the standard L1 noise-prediction
loss on the diffusion target,
$\mathcal{L}_\text{percep}$ is the VGG16-feature perceptual loss
\cite{simonyan2014vgg}, and $\mathcal{L}_\text{topo}$ is the proposed
edge-aware topology loss
%
\begin{equation}
\mathcal{L}_\text{topo} = \mathcal{L}_\text{CE}(\hat S,S) + \lambda_e \mathcal{L}_\text{CE}(\nabla\hat S,\nabla S) + \lambda_d \mathcal{L}_\text{Dice}(\hat S,S),
\end{equation}
%
where $\hat S$ is the segmentation logits, $S$ is the 7T-derived
ground truth, $\nabla$ is a 3D Sobel gradient operator applied per class,
and $\lambda_e=0.5$, $\lambda_d=0.5$. The Sobel-gradient term penalises
ill-defined boundaries that drive the connected-component
fragmentation observed in early epochs (\S\ref{sec:res-ablation}), while
the Dice component encourages volumetric agreement. Training uses Adam
\cite{kingma2014adam} with learning rate $10^{-4}$, batch size 8, weight
decay $10^{-4}$, gradient clipping at norm 1.0, cosine diffusion
schedule with $T=200$ timesteps, and EMA tracking from iteration 2{,}000
with decay 0.9999.

\subsection{Evaluation metrics}
\label{sec:metrics}

Synthesis quality is reported with structural similarity index measure
(SSIM, \cite{wang2004image}), peak signal-to-noise ratio (PSNR), Dice
coefficient (foreground and per tissue class), 95\,th-percentile
Hausdorff distance (HD95) for surface fidelity, and average symmetric
surface distance (ASSD). Anatomical preservation is reported through
per-tissue connected-component counts (CC), the fraction of voxels
captured by the largest connected component of the predicted class
(\emph{largest-CC fraction}), and per-region GM coverage inside the
ground-truth aparc+aseg ROIs. For ADNI deployment, AD vs CN tissue-
volume Cohen's $d$, group means, and a NIA-AA-style composite
biomarker score \cite{jack2018niaaa} provide quantitative biomarker
analysis.

% =============================================================================
\section{Results}
\label{sec:results}
% =============================================================================

\subsection{Held-out paired evaluation (sub-06)}
\label{sec:res-paired}

Table~\ref{tab:paired-quality} reports synthesis quality on the
held-out subject. The 105\,k-iteration EMA checkpoint was selected by
candidate screening on PSNR followed by full-volume SSIM, Dice, HD95
and qualitative review.

\begin{table}[!t]
\caption{Synthesis quality on held-out subject sub-06 at the
selected 105\,k EMA checkpoint.}
\label{tab:paired-quality}
\centering
\small
\begin{tabular}{lc}
\toprule
Metric & Value \\
\midrule
SSIM (whole volume)              & 0.8991 \\
PSNR (whole volume, dB)          & 20.00 \\
Dice (foreground)                & 0.7682 \\
HD95 (foreground, mm)            & 3.31 \\
Brain-mask Dice                  & 0.9436 \\
Brain-mask HD95 (mm)             & 2.05 \\
\bottomrule
\end{tabular}
\end{table}

\subsection{Per-tissue and topology metrics}
\label{sec:res-tissue}

Table~\ref{tab:tissue-dice} summarises tissue-class fidelity. Grey
matter and white matter Dice exceed 0.85, whereas CSF --- the thinnest
tissue class with strongest boundary sensitivity --- reaches Dice
0.7616 (compared with 0.4785 for the no-topology variant; see
\S\ref{sec:res-ablation}).

\begin{table}[!t]
\caption{Per-tissue Dice and connected-component analysis on sub-06.
Largest-CC fraction is the proportion of predicted-class voxels
captured by the largest connected component (higher is better).}
\label{tab:tissue-dice}
\centering
\small
\begin{tabular}{lcccc}
\toprule
Tissue & Dice & Vol.\ diff. (\%) & CC pred & Largest-CC \\
\midrule
CSF        & 0.7616 & $+36.5$  &   126   & 0.641 \\
GM         & 0.8501 & $+12.7$  & 1{,}658 & 0.983 \\
WM         & 0.8963 & $+4.2$   & 1{,}038 & 0.965 \\
Brain      & 0.9436 & ---      & ---     & ---   \\
\bottomrule
\end{tabular}
\end{table}

\subsection{Topology-loss ablation}
\label{sec:res-ablation}

Table~\ref{tab:ablation} compares the full TopoBrain model against an
ablated variant trained without the topology loss. Whole-brain HD95
drops from 13.34\,mm (no-topology) to 2.05\,mm (with topology), a
6.5$\times$ reduction. Component fragmentation falls across all
tissue classes, including a 7$\times$ reduction in CSF components and
a 1.6$\times$ reduction in WM components.

\begin{table}[!t]
\caption{Effect of the proposed edge-aware topology loss
($\lambda=0.2$) on the held-out subject. The no-topology variant is
the same architecture and schedule trained with
$\mathcal{L}_\text{total}=\mathcal{L}_\text{pixel}+0.25\mathcal{L}_\text{percep}$
through 50\,k iterations; the topology variant is the selected 105\,k
EMA checkpoint.}
\label{tab:ablation}
\centering
\small
\begin{tabular}{lcc}
\toprule
Metric & No-topology & TopoBrain \\
\midrule
Brain-mask Dice            & 0.9054 & \textbf{0.9436} \\
Brain-mask HD95 (mm)       & 13.34  & \textbf{2.05}   \\
CSF Dice                   & 0.4785 & \textbf{0.7616} \\
GM Dice                    & 0.7949 & \textbf{0.8501} \\
WM Dice                    & 0.8512 & \textbf{0.8963} \\
CSF components             & 1{,}008 & \textbf{126}   \\
GM components              & 1{,}889 & \textbf{1{,}658}\\
WM components              & 1{,}645 & \textbf{1{,}038}\\
\bottomrule
\end{tabular}
\end{table}

\subsection{Subcortical region preservation (AD-relevant ROIs)}
\label{sec:res-roi}

Table~\ref{tab:roi} reports subcortical preservation in the eight
AD-relevant ROIs defined by the ground-truth FreeSurfer parcellation.
Bilateral hippocampi, amygdalae, thalami and caudates are recovered as
essentially single connected components with largest-CC fraction
$\geq$\,99.8\,\% and GM coverage $\geq$\,81.8\,\%. This is the
prerequisite for downstream subcortical volumetry on synthesised
volumes; the systematic $-5$\,\% to $-18$\,\% under-segmentation is
uniform across regions and acceptable provided it is uniform across
clinical subgroups (verified for AD/CN in \S\ref{sec:res-adni}).

\begin{table}[!t]
\caption{Subcortical region preservation on the held-out subject.
``Coverage'' is the fraction of voxels in the GT region also predicted
as GM in the synthesised volume.}
\label{tab:roi}
\centering
\small
\begin{tabular}{lccc}
\toprule
Region & CC & Largest-CC & Coverage \\
\midrule
Hippocampus L  & 6 & 0.998 & 0.852 \\
Hippocampus R  & 2 & 1.000 & 0.922 \\
Amygdala L     & 1 & 1.000 & 0.818 \\
Amygdala R     & 1 & 1.000 & 0.954 \\
Thalamus L     & 1 & 1.000 & 0.991 \\
Thalamus R     & 1 & 1.000 & 0.987 \\
Caudate L      & 1 & 1.000 & 0.943 \\
Caudate R      & 1 & 1.000 & 0.954 \\
\bottomrule
\end{tabular}
\end{table}

\subsection{Out-of-distribution ADNI deployment}
\label{sec:res-adni}

Figure~\ref{fig:cohort} shows the synthesised mid-coronal slice for
each of the 30 ADNI baseline subjects. All 30 subjects produce
anatomically coherent outputs with no catastrophic artefacts; cohort-
level GM topology preservation reaches largest-CC fraction $0.847
\pm 0.030$, white-matter $0.990 \pm 0.006$ --- statistically equivalent
to the training-cohort baseline (sub-06: 0.85), demonstrating that
anatomical preservation generalises to the older, multi-scanner
out-of-distribution cohort.

\begin{figure*}[!t]
\centering
\includegraphics[width=0.9\textwidth]{cohort_overview.png}
\caption{Synthesised 7T outputs for the 30 ADNI baseline 3T subjects
(15 AD, 15 CN), mid-coronal slice. All subjects produce anatomically
coherent volumes; group-level mean grey-matter largest-connected-
component fraction is $0.847 \pm 0.030$, white-matter $0.990 \pm 0.006$.}
\label{fig:cohort}
\end{figure*}

Table~\ref{tab:adni-volumes} reports group-mean tissue volumes
extracted from the model's predicted segmentation. All six canonical
biomarker directions match the AD literature
\cite{frisoni2010clinical,jack2018niaaa}: AD shows lower brain, GM, WM,
and GM fraction; AD shows higher CSF and CSF/GM atrophy index. The
strongest single signal is the CSF/GM atrophy index with Cohen's
$d=+0.681$ (medium effect), followed by CSF volume ($d=+0.561$) and
GM volume ($d=-0.507$).

\begin{table}[!t]
\caption{Group-level tissue volumes on the 30-subject ADNI cohort
(15 AD, 15 CN). All directions match the AD biomarker literature.}
\label{tab:adni-volumes}
\centering
\small
\begin{tabular}{lccc}
\toprule
Volume & AD mean & CN mean & Cohen's $d$ \\
\midrule
Brain (cm$^3$)           & 1055.3 & 1099.2 & $-0.45$ \\
GM (cm$^3$)              &  599.5 &  629.1 & $-0.51$ \\
WM (cm$^3$)              &  419.4 &  440.4 & $-0.34$ \\
CSF (cm$^3$)             &   36.3 &   29.6 & $+0.56$ \\
GM fraction              & 0.5708 & 0.5742 & $-0.18$ \\
\textbf{CSF/GM ratio}    & \textbf{0.0607} & \textbf{0.0471} & $\mathbf{+0.68}$ \\
\bottomrule
\end{tabular}
\end{table}

\subsection{AD biomarker analysis on ADNI}
\label{sec:res-biomarker}

A NIA-AA-style composite z-score \cite{jack2018niaaa} on the four-
biomarker panel \{GM fraction, GM volume, CSF/GM ratio, CSF volume\}
yields exploratory area-under-curve $=$\,0.653 with sensitivity
0.400 and specificity 0.933 at the Youden-J operating point
(Fig.~\ref{fig:roc}). This is below the 3T-baseline composite (AUC
$=$\,0.818) computed from naive K-means tissue thresholding of the
input 3T volume on the same cohort, indicating that at the four-class
tissue granularity the AD signal is already captured by simple
3T-only segmentation. The synth-derived score remains above chance
and the per-feature Cohen's $d$ values match the published-literature
direction in 6/6 cases. The mechanistic argument for the synthesis is
therefore not coarse tissue-volume improvement; rather, the
hypothesised advantage of 7T-supervised segmentation manifests at
finer parcellation (hippocampal sub-fields, entorhinal cortex
thickness, regional cortical volumes), which is the immediate next
study (\S\ref{sec:future}).

\begin{figure}[!t]
\centering
\includegraphics[width=\columnwidth]{ad_risk_distribution.png}
\caption{Composite NIA-AA-style AD-risk z-score (synth-derived
features) on the 30-subject ADNI cohort. Positive values indicate
greater AD-likelihood. Group means are shown as solid horizontal lines
and the Youden-J optimal threshold as a dashed line.}
\label{fig:roc}
\end{figure}

% =============================================================================
\section{Discussion}
\label{sec:discussion}
% =============================================================================

\textbf{Anatomical fidelity vs.~perceptual sharpness.} The synthesised
volumes preserve subcortical anatomy at the millimetre scale
(Table~\ref{tab:roi}, Fig.~\ref{fig:compare}) but exhibit reduced
high-frequency cortical texture relative to the input 3T volumes. This
is consistent with the perception-distortion trade-off
\cite{blau2018perception} for MSE-trained generative models on small
training cohorts: pixel-similarity objectives recover the conditional
mean, which is mathematically smoother than any individual sample.
Notably, this smoothing does not impair the downstream AD-detection
task, because AD-relevant biomarkers --- hippocampal volume, CSF/GM
atrophy index, ventricular enlargement --- are low-frequency
structural features that survive smoothing.

\begin{figure}[!t]
\centering
\includegraphics[width=\columnwidth]{006_S_4153_compare.png}\par
\vspace{2pt}
\includegraphics[width=\columnwidth]{002_S_4213_compare.png}
\caption{Three-orthogonal-plane comparisons for one AD subject (top:
006\_S\_4153, age 79.6) and one CN subject (bottom: 002\_S\_4213,
age 78.7). Top row in each panel: preprocessed 3T input; bottom row:
synthesised 7T output. Cortical-ribbon texture is reduced (perception-
distortion trade-off) but subcortical anatomy --- hippocampus,
amygdala, thalamus, caudate --- and ventricular morphology are
preserved.}
\label{fig:compare}
\end{figure}

\textbf{The role of the topology loss.} The ablation in
Table~\ref{tab:ablation} shows that the topology term contributes
substantially to whole-brain boundary fidelity (HD95 6.5$\times$
lower) and reduces per-tissue component fragmentation across CSF, GM
and WM. The current implementation uses Sobel-edge cross-entropy and
boundary Dice; we explicitly note that this is an \emph{edge-aware
topology approximation}, not a persistent-homology loss in the strict
sense. Replacement with persistent-homology-based losses
\cite{hu2019topology,clough2020topo} or centerline-Dice
\cite{shit2021cldice} is identified as worthwhile follow-up.

\textbf{Granularity of the AD comparison.} The 30-subject ADNI
deployment yields six-of-six biomarker-direction agreement with the
literature but a composite-score AUC of 0.653 versus 0.818 for the
3T baseline. We interpret this as a granularity issue rather than a
model failure: at the four-class tissue level, the AD signal is
already strong in 3T, and naive thresholding captures it. The
synthesis pipeline's mechanistic advantage is access to 7T-supervised
\emph{tissue boundaries}, which is most likely to manifest at finer
parcellation --- hippocampal subfields (e.g., CA1 \cite{wisse2014hippo}),
entorhinal-cortex thickness, regional cortical thickness via Desikan-
Killiany atlases. The full FastSurfer/FreeSurfer parcellation
evaluation on the full 200-subject age/sex-matched cohort is the
natural next study; the pipeline for this evaluation is implemented
and self-validated against synthetic data.

\textbf{Clinical relevance.} The proposed synthesis pipeline addresses
a real clinical accessibility gap: $\sim$99\,\% of clinical AD imaging
worldwide is at 3T, whereas the published 7T biomarker advantages
\cite{wisse2014hippo,deflores2020hippocampal} are anchored on
$<$1\,\% of available scanner installations. A synthesis model
producing 7T-quality contrast on the patient's own 3T voxel grid, at
3.1\,min/subject inference, lowers the deployment threshold and
enables prospective validation studies on otherwise-7T-inaccessible
cohorts.

% =============================================================================
\section{Limitations and Future Work}
\label{sec:future}
% =============================================================================

\textbf{Training cohort size.} Training on $N=10$ paired subjects is
the principal limitation; high-frequency cortical detail is the
expected casualty of this small training set
\cite{blau2018perception}. Larger paired 3T/7T cohorts (e.g., the
publicly distributed BICCN, MGH 7T initiatives) would directly
mitigate this and are an immediate priority.

\textbf{ADNI cohort statistical power.} The 30-subject deployment
provides effect-size estimates and direction-of-effect evidence; it is
not powered for inferential statistical claims about AD
discriminability. The full pipeline --- including paired DeLong's
testing for AUC comparison \cite{sun2014delong} and intra-class
correlation analysis \cite{shrout1979icc} for volumetry agreement
--- is implemented; running it on the matched 200-subject cohort with
FastSurfer-derived parcellated features is the natural next study.

\textbf{Topology loss formulation.} The edge-aware Sobel-Dice
combination used here approximates topology preservation through
boundary supervision. Persistent-homology losses
\cite{hu2019topology,clough2020topo} or centerline-Dice variants
\cite{shit2021cldice} are theoretically stronger candidates and merit
direct comparison against our implementation.

\textbf{Single-modality input.} Only T1w is processed. Multi-modal
synthesis (T1 + T2 + FLAIR) is expected to improve robustness and
discrimination of pathology-adjacent tissue.

\textbf{Inference latency.} The 200-step DDPM sampler used for
training-time fidelity is reduced to 50 DDIM steps for inference
without measurable degradation; further acceleration via consistency
distillation is a natural extension.

% =============================================================================
\section{Conclusion}
\label{sec:conclusion}
% =============================================================================

We have presented \textbf{TopoBrain}, an anatomy-aware conditional
diffusion framework for 3T-to-7T MRI synthesis that introduces an
auxiliary tissue-segmentation head and an edge-aware topology loss to
preserve subcortical anatomy relevant to early Alzheimer's disease
biomarkers. On a held-out paired evaluation, TopoBrain achieves
SSIM\,$=$\,0.8991, PSNR\,$=$\,20.00\,dB, brain-mask Dice 0.9436 and
HD95\,$=$\,2.05\,mm; the topology loss reduces brain HD95 by
6.5$\times$ relative to a no-topology ablation, and bilateral
hippocampi, amygdalae, thalami and caudates are recovered with
largest-component fraction $\geq$\,0.998. Out-of-distribution
deployment on 30 ADNI baseline 3T subjects (15\,AD, 15\,CN) produces
anatomically coherent outputs whose tissue-volume biomarkers recover
six-of-six expected AD directions with a CSF/GM atrophy-index Cohen's
$d=+0.681$. We have implemented and self-validated an end-to-end
downstream pipeline; full empirical AD-classification on the matched
200-subject ADNI cohort is the immediate next study, requiring
large-scale FastSurfer-based parcellation compute beyond the present
scope. All code, configurations and analysis scripts are openly
available\footnote{\url{https://github.com/prabeshx12/Topo-Brain}}.

% =============================================================================
\section*{Acknowledgements}
% =============================================================================

The authors thank the Department of Electronics and Computer
Engineering, Pulchowk Campus, Tribhuvan University, and Dr.\ Nanda
Bikram Adhikari for academic supervision; E.K.\ Solutions for
providing access to NVIDIA A100 GPU resources; and the Alzheimer's
Disease Neuroimaging Initiative for distributing the 3T MPRAGE data
used in the external-cohort evaluation.

% =============================================================================
\section*{Conflict of interest}
% =============================================================================

The authors declare no conflict of interest.

% =============================================================================
\section*{Data availability statement}
% =============================================================================

ADNI data used in this study were obtained from the public Alzheimer's
Disease Neuroimaging Initiative archive
(\url{https://adni.loni.usc.edu}) under the standard ADNI Data Use
Agreement. The paired 3T/7T training cohort
\cite{chen2023unc} is distributed by its original maintainers; we do
not redistribute. All analysis code, including cohort-assembly
tooling, synthesis scripts, parcellation and classifier-comparison
pipelines, and self-tests, is available at
\url{https://github.com/prabeshx12/Topo-Brain}.

% =============================================================================
\section*{Author contributions}
% =============================================================================

\textbf{P.B.}, \textbf{P.A.}, and \textbf{S.B.} jointly designed the
study, developed the pipeline, and analysed the results.
\textbf{N.B.A.} supervised the project and provided methodological
guidance. All authors wrote and reviewed the manuscript.

% =============================================================================
%  References
% =============================================================================
\bibliographystyle{IEEEtran}
\bibliography{references}

\end{document}
